\input{../tex-inputs/latex-leseansicht-vorspann}

\section[Arthur Schnitzler an Gustav Schwarzkopf, 12. 8. 1926]{L04462 Arthur Schnitzler an Gustav Schwarzkopf, 12. 8. 1926}
\nopagebreak\mylabel{L04462v}
\rehead{ }\normalsize\beginnumbering\briefempfaengerindex{Schwarzkopf, Gustav@\textsc{Schwarzkopf, Gustav}!zzzSchnitzler, Arthur@\emph{von Arthur Schnitzler}!1926-08-121@{12. 8. 1926}|(be}
\toendnotes[C]{\smallbreak\pagebreak[2]}
\correspDesc{Versand  durch Arthur Schnitzler am 12. 8. 1926 in Bern
\newline{}Übermittlung  am 12. 8. 1926 in Bern
\newline{}Erhalt  durch Gustav Schwarzkopf im Zeitraum [13. 8. 1926 – 17. 8. 1926?] in Wien}\toendnotes[C]{\smallbreak}
\Standort{DLA, A:Schnitzler, HS.1985.1.1897.}
\physDesc{Bildpostkarte, 390 Zeichen
\newline{}Handschrift: Bleistift, lateinische Kurrent
\newline{}Versand: 1) Stempel: »\nobreak{}\oindex{Bern@\textbf{Bern}|pwk}Bern 1, 12–VIII 192\textcolor{gray}{6}, 23–24, Briefversand\nobreak{}«.   2) Stempel: »\nobreak{}Ferien in der Schweiz\oindex{Schweiz@\textbf{Schweiz}|pw}\nobreak{}«. }\toendnotes[C]{\smallbreak}\pstart{}{\pb}Hrn Gustav Schwarzkopf\pend{}\pstart{}Wien I\oindex{I., Innere Stadt@\textbf{I., Innere Stadt}|pw}\pend{}\pstart{}Tiefer Graben \substVorne{}\textsuperscript{2}\substDazwischen{}1\substHinten{}7\oindex{Wien@\textbf{Wien}!I., Innere Stadt@\textbf{I., Innere Stadt}!Tiefer Graben 17@\textbf{Tiefer Graben 17}|pw}\pend{}{\bigskip}
\pstart
           \noindent{}{\pb}\textcolor{gray}{\textbf{Bern\oindex{Bern@\textbf{Bern}|pw}. Südseite}}\pend
           \vspace{1em}
\pstart
           \raggedleft{}{\pb}12. 8.\pend
           \vspace{0.5em}
\pstart
           Vielen Dank für die Nachrichten, lieber Guſtav, die ja im Ganzen als abzureden
      gelten dürfen. – Ich bin seit gestern
      hier, und weiſs noch nicht, ob ich 
      zunächst an den Thuner\oindex{Thunersee@\textbf{Thunersee}|pw} oder an
      den Genfer See\oindex{Genfer See@\textbf{Genfer See}|pw} gehen werde. Wetter
      war meist nicht gut. Alles dankt
      herzlichst für Ihre lieben Grüße.
      Heini\pwindex{Schnitzler, Heinrich 9.\,8.\,1902 Hinterbrühl – 12.\,7.\,1982 Wien@\textsc{Schnitzler, Heinrich} (9.\,8.\,1902 Hinterbrühl – 12.\,7.\,1982 Wien)|pw} ist eben nach Berlin\oindex{Berlin@\textbf{Berlin}|pw}, Lili\pwindex{Cappellini, Lili 13.\,9.\,1909 Wien – 26.\,7.\,1928 Venedig@\textsc{Cappellini, Lili} (13.\,9.\,1909 Wien – 26.\,7.\,1928 Venedig)|pw}
      mit O\pwindex{Schnitzler, Olga 17.\,1.\,1882 Wien – 13.\,1.\,1970 Lugano@\textsc{Schnitzler, Olga} (17.\,1.\,1882 Wien – 13.\,1.\,1970 Lugano)|pw} in Venedig\oindex{Venedig@\textbf{Venedig}|pw}.\pend
           \pstart Auf Wiedersehen \label{T_L04462-1v}\edtext{Ihr \spacefill\mbox{A}}{\lemma{\textnormal{\emph{Ihr A}}}\Cendnote{\textnormal{entlang des rechten Randes geschrieben}}}\label{T_L04462-1}\pend{}\selectlanguage{ngerman}\endnumbering\briefempfaengerindex{Schwarzkopf, Gustav@\textsc{Schwarzkopf, Gustav}!zzzSchnitzler, Arthur@\emph{von Arthur Schnitzler}!1926-08-121@{12. 8. 1926}|)be}\mylabel{L04462h}
\begin{anhang}
\end{anhang}\newcommand{\dateiname}{L04462}\newcommand{\titel}{Arthur Schnitzler an Gustav Schwarzkopf, 12. 8. 1926}\newcommand{\editorInnen}{Herausgegeben von Jahnke, SelmaMüller, Martin Anton}\input{../tex-inputs/latex-leseansicht-abspann}
